\chapter{ПРЕДВАРИТЕЛЬНОЕ ТЕСТИРОВАНИЕ НА ИСКУССТВЕННЫХ ДАННЫХ И СБОР ЭТАЛОННЫХ ПОКАЗАТЕЛЕЙ}
\label{cha:test}

\section{Методика тестирования}

Предварительное тестирование прототипа проводилось на \emph{искусственных
(синтетических)} данных --- сценах с заранее известными свойствами, позволяющих
оценить корректность, устойчивость и производительность без привязки к
конкретному прикладному сценарию. Тестирование преследовало три цели:
\begin{enumerate}
  \item проверка функциональной корректности модулей (модульные тесты);
  \item проверка \emph{детерминизма} --- воспроизводимости результата при
        повторных запусках на одной и той же сборке и аппаратной платформе;
  \item сбор \emph{эталонных показателей} производительности, используемых как
        опорные значения на последующих этапах.
\end{enumerate}

\section{Модульное тестирование}

Физическое ядро \texttt{dyphur} покрыто набором модульных тестов (фреймворк
Catch2): вычислительный слой (буферы, параллельные примитивы, редукция,
сортировка), интегратор, широкая и узкая фазы обнаружения столкновений, решатель
XPBD, сочленения и трассировка лучей. Тесты снабжены метками (\texttt{smoke},
\texttt{determinism}), что позволяет выборочно запускать быстрые проверки в
процессе разработки и полный набор --- в непрерывной интеграции.

\section{Эталонный сценарий и показатели}

В качестве эталонного нагрузочного сценария используется укладка \textbf{1024}
твёрдых тел (боксов) на статическую опорную плоскость; моделируется 5~секунд
с частотой 60~Гц. Для каждого прогона фиксируются:
\begin{itemize}
  \item средняя частота кадров симуляции (FPS);
  \item коэффициент реального времени (отношение модельного времени к
        затраченному);
  \item хеш итогового состояния --- для контроля детерминизма.
\end{itemize}

Результаты на двух аппаратных конфигурациях приведены в
таблице~\ref{tab:bench}. Совпадение хешей между повторными запусками на одной
конфигурации подтверждает детерминизм; различие хешей CPU и GPU ожидаемо и
обусловлено различным порядком операций с плавающей точкой (см.
раздел~\ref{cha:gpu}).

\begin{table}[ht]
  \centering
  \caption{Эталонные показатели для сценария 1024 тел (5~с, 60~Гц)}
  \label{tab:bench}
  \small
  \begin{tabular}{p{0.30\textwidth} p{0.16\textwidth} r r}
    \toprule
    Аппаратура & Бэкенд & FPS & Реальное время \\
    \midrule
    Intel Xeon E5-2667 v4 & OpenMP (CPU) & 254 & 4{,}24$\times$ \\
    NVIDIA RTX 3060       & CUDA (GPU)   & 83  & 1{,}39$\times$ \\
    \bottomrule
  \end{tabular}
\end{table}

\screenshot{0.85}{Сбор эталонных показателей на синтетическом сценарии}{Снимок
экрана: вывод нагрузочного теста (1024 тел) --- значения FPS, коэффициента
реального времени и хеша состояния; либо график зависимости FPS от числа
тел.}{benchmark}

На текущем этапе (фаза~1) на обеих платформах достигнута симуляция 1024~тел в
реальном времени с побитово воспроизводимым результатом. Полученные значения
зафиксированы как эталонные: относительно них на последующих этапах будет
оцениваться эффект от оптимизации GPU-конвейера и масштабирования числа тел.

\section{Результаты}

Сформирована методика предварительного тестирования на искусственных данных,
реализован набор модульных и сквозных тестов, подтверждены функциональная
корректность и детерминизм ядра, собраны и зафиксированы эталонные показатели
производительности для опорного нагрузочного сценария.
